\documentclass[12pt,a4paper]{article}

\usepackage[utf8]{inputenc}
\usepackage[T1]{fontenc}
\usepackage[french]{babel}
\usepackage{amsmath,amssymb,amsthm}
\usepackage{cmbright} % whole document sans-serif, matching sans math
\usepackage{exscale}
\usepackage[hidelinks]{hyperref}
\usepackage{cleveref}
\usepackage{graphicx}
\usepackage{tikz}
\usepackage{caption}
\usepackage{standalone}
\usetikzlibrary{decorations.pathreplacing}
\usepackage{xcolor}
\definecolor{accent}{HTML}{244470} % not exactly navy blue, but better imo
\usepackage{geometry}
\geometry{hmargin=2cm,vmargin=2cm}
\usepackage{enumitem}

\usepackage{titlesec}
% Style the sections with the navy color
\titleformat{\section}
{\color{accent}\Large\bfseries}
{\thesection.}
{0.8em}
{}
\titleformat{\subsection}
{\color{accent}\large\bfseries}
{\thesubsection}
{0.7em}
{}

\usepackage[many]{tcolorbox}
% Create a box that is white, but has a 3pt line on the left
\tcbset{
% Use for Theorems and Definitions
navybar/.style={
enhanced,
colback=white,
colframe=accent,
boxrule=0.6pt,
leftrule=3pt,
arc=0pt,
outer arc=0pt,
left=8pt, right=8pt, top=6pt, bottom=6pt
},
% Use for properties / less important results
graybar/.style={
enhanced,
colback=white,
colframe=gray!70,
boxrule=0.6pt,
leftrule=3pt,
arc=0pt,
outer arc=0pt,
left=8pt, right=8pt, top=6pt, bottom=6pt
},
}

% Theorems and definitions
\theoremstyle{plain}
\newtheorem{theorem}{Théorème}[section]
\newtheorem{property}[theorem]{Propriété}

\theoremstyle{definition}
\newtheorem{definition}[theorem]{Définition}
\newtheorem{example}[theorem]{Exemple}

\newtheorem{exercise}{Exercice}

\theoremstyle{remark}
\newtheorem*{remark}{Remarque}
\newtheorem*{recall}{Rappel}

\begin{document}

\noindent
\begin{minipage}[c]{0.62\textwidth}
  {\color{accent}\LARGE\bfseries Généralités sur les fonctions\par}
  \vspace{0.1cm}
  {\large 3\textsuperscript{ème} année secondaire\par}
  \vspace{0.7cm}
  {\bfseries \href{https://www.linkedin.com/in/kefrov/}{Fares Manai}\par}
  \vspace{2pt}
  {\small\href{mailto:kefrov630@gmail.com}{kefrov630@gmail.com}\par}
\end{minipage}
\hfill
\begin{minipage}[c]{0.33\textwidth}
  \centering
  \includegraphics[width=\linewidth]{fx.pdf}
\end{minipage}
\par\vspace{0.4cm}
\noindent{\color{accent}\rule{\linewidth}{0.9pt}}
\par\vspace{0.5cm}

\tableofcontents

\newpage
\phantomsection
\section*{Introduction}
\addcontentsline{toc}{section}{Introduction}

\newpage
\section{Rappels}

\begin{tcolorbox}[navybar]
  \begin{definition}
    Given two sets $A$ and $B$, a function from $A$ to $B$ is a rule or mapping
    that takes each element $x \in A$ and associates with it a single element of $B$.
    Given an element $x \in A$, the expression $f(x)$ is used to represent the element
    of $B$ associated with $x$ by $f$. In this case, we write
    \[
    \begin{aligned}
      f : A &\rightarrow B \\
      x &\mapsto f(x).
    \end{aligned} 
    \]    
    The set $A$ is called the domain of $f$, and $B$ the codomain. The range of $f$ 
    refers to the set of all possible outputs of $f$ given by 
    \[
    f(A) = \{y \in B \mid y = f(x) \text{ for some } x \in A\}. 
    \]
  \end{definition}
\end{tcolorbox}

The first thing we emphasize about functions is that each input has a unique output.
For example, stating for some $x$ that both $f(x) = 1$ and $f(x) = 2$ hold makes no 
sense based on the definition of the function. 

Two functions with the same expression
but different domains are considered different. In other words, a function is defined 
by its domain and its rule or expression. 

\begin{remark}
  Be careful not to confuse the function $f$ with its expression $f(x)$. Note also
  that the range of $f$ is not necessarily its codomain, the range is a subset 
  of the codomain. 
\end{remark}

Visualizing functions defined on $\mathbb{R}$ graphically could be very useful durant 
leur etude. Hence the following definition.
\begin{tcolorbox}[navybar]
  \begin{definition}
    Given a plane $(O, \vec{i}, \vec{j})$, the courbe of a function $f : D_f \rightarrow
    \mathbb{R}$ is the set of points 
    \[
    \{(x, f(x)) \mid x \in D_f\}.
    \]
  \end{definition} 
\end{tcolorbox}

\begin{example}
  There can't exist a function that to each $x$ associates a $y$ based on the rule
$y^2 = x$, as that rule associates two values $y$ and $-y$ to each $x$. 
  In Figure~\ref{fig:inverse_x_sq}, we illustrate the set of points
$(x,y)$ that satisfy such a rule. Notice from the vertical line with equation
$x=4$ that two values of $y$, $-2$ and $2$, correspond to $x$.
  
  If we were to consider $y$ as the input and $x$ as the output, then it is now
  possible to define a function, which would simply be $y \mapsto y^2$.
\end{example}

\begin{example}
  Consider the functions $f : \mathbb{R}_- \rightarrow \mathbb{R}$ and 
  $g : \mathbb{R}_+ \rightarrow \mathbb{R}$ such that $f$ and $g$ share the 
  same expression $x^2$. Figure~\ref{fig:x_sq} makes it apparent that these 
  functions are not the same through their graphs $C_f$ and $C_g$.
\end{example}

We will now go over a set of useful functions that will be frequently appearing, 
starting with the polynomial.

\begin{tcolorbox}[navybar]
  \begin{definition}
    A function $p : \mathbb{R} \rightarrow \mathbb{R}$ is said to be a polynomial
    if its expression for all $x \in \mathbb{R}$ is of the form
    \[
    p(x) = a_0 + a_1 x + a_2 x^2 + \cdots + a_m z^m,
    \] 
    where $a_0, \dots, a_m \in \mathbb{R}$.
  \end{definition} 
\end{tcolorbox}

In the 2eme annee secondaire, we mostly went over polynomials where $m = 2$, these
polynomials are said to be of second degree.

\noindent
\begin{minipage}[t]{0.38\linewidth}
  \vspace{-10pt}
  \captionsetup{hypcap=false}
  
  \begin{minipage}[c][6cm][c]{\linewidth}
    \centering
    \includestandalone[width=\linewidth]{figures/inverse_x_sq}
  \end{minipage}
  
  \captionof{figure}{$\{(x,y) \in \mathbb{R}^2 \mid y^2=x\}$.}
  \label{fig:inverse_x_sq}
\end{minipage}\hfill
\begin{minipage}[t]{0.38\linewidth}
  \vspace{-10pt}
  \captionsetup{hypcap=false}
  
  \begin{minipage}[c][6cm][c]{\linewidth}
    \centering
    \includestandalone[width=\linewidth]{figures/x_sq}
  \end{minipage}
  
  \captionof{figure}{$C_f$ et $C_g$.}
  \label{fig:x_sq}
\end{minipage}
\par

\begin{tcolorbox}[navybar]
  \begin{definition}
    A polynomial $p : \mathbb{R} \rightarrow \mathbb{R}$ is said to be of degree $m$
    if for all $x \in \mathbb{R}$ its expression is of the form 
    \[
    p(x) = a_0 + a_1 x + a_2 x^2 + \cdots + a_m z^m,
    \] 
    where $a_0, \dots, a_m \in \mathbb{R}$ and $a_m \neq 0$.
  \end{definition} 
\end{tcolorbox}

We were also particularly interested in the roots of degree $2$ polynomials. 

\begin{tcolorbox}[navybar]
  \begin{definition}
    The roots of a polynomial $p$ are the reals $x$ such that $p(x) = 0$.
  \end{definition} 
\end{tcolorbox}

\noindent
\begin{minipage}[t]{0.48\linewidth}
  \vspace{0pt}
  \begin{example}
    Considérons l'équation du second degré
    \[ 
      x^2 - 5x + 6 = 0. 
    \]
    Graphically, as shown in Figure~\ref{fig:poly_roots}, 
    the roots are $2$ and $3$, which you can also easily verify by calculating $p(2)$
    and $p(3)$. 
    
    One question arises however, how can we \textit{reliably} find the roots?
    i.e., without resorting to the courbe of $p$.
  \end{example}
\end{minipage}\hfill
\begin{minipage}[t]{0.48\linewidth}
  \vspace{0pt}
  \captionsetup{hypcap=false}
  
  \begin{minipage}[c][6cm][c]{\linewidth}
    \centering
    \includestandalone[width=\linewidth]{figures/poly_roots}
  \end{minipage}
  
  \captionof{figure}{Roots of $x^2 - 5x + 6$.}
  \label{fig:poly_roots}
\end{minipage}
\par

Consider the equation \(ax^2 + bx + c = 0\) where \(a \neq 0\). Graphically, we see that
this equation could have either 0, 1 or 2 solutions, which would be the roots of
the polynomial \(ax^2 + bx + c\). So a natural question arises: under what conditions
do we have each of these cases, and what are the forms of such roots?

Notice that there are two terms containing \(x\), being \(ax^2\) and \(bx\). The algebraic
intuition here is to possibly reformulate the expression so that there is only 
one instance of \(x\) instead of two. Observe the following remarkable product
\[
(x + A)^2 = x^2 + A^2 + 2Ax.
\]
The expression on the left has only one instance of \(x\), whereas on the right we have two.
Trying to exploit this for our equation yields
\[
\begin{aligned}
  x^2 + \frac{b}{a}x + \frac{c}{a} &= 0 \quad \text{(dividing by } a \neq 0 \text{)} \\
  x^2 + 2\left(\frac{b}{2a}\right)x + \frac{c}{a} &= 0 \quad \text{(rewriting the middle term)} \\
  x^2 + 2\left(\frac{b}{2a}\right)x + \left(\frac{b}{2a}\right)^2 - \left(\frac{b}{2a}\right)^2 + \frac{c}{a} &= 0 \quad \text{(adding and subtracting the square)} \\
  \left(x + \frac{b}{2a}\right)^2 - \frac{b^2}{4a^2} + \frac{c}{a} &= 0 \quad \text{(factoring the perfect square)} \\
  \left(x + \frac{b}{2a}\right)^2 - \frac{b^2 - 4ac}{4a^2} &= 0 \quad \text{(finding a common denominator)} \\
  \left(x + \frac{b}{2a}\right)^2 &= \frac{b^2 - 4ac}{4a^2}
\end{aligned}
\]

We now have exactly what we wanted: a single instance of \(x\) on the left side. 
The term \(b^2 - 4ac\) is so important that we give it a special name: the discriminant, denoted by \(\Delta\).
\[
\left(x + \frac{b}{2a}\right)^2 = \frac{\Delta}{4a^2}
\]

Because the left side is a square, it must be non-negative. Since the denominator \(4a^2\) is strictly positive, the number of real solutions depends entirely on the sign of \(\Delta\). Hence we get the following 3 cases:

\begin{itemize}
  \item If \(\Delta < 0\): The right side is negative. Since a square cannot be negative, there are 0 solutions.
  \item If \(\Delta = 0\): The right side is zero. Thus \(x + \frac{b}{2a} = 0\), yielding exactly 1 solution (a double root): 
  \[
  x = -\frac{b}{2a}
  \]
  \item If \(\Delta > 0\): We can take the square root of both sides, yielding 2 solutions:
  \[
  x + \frac{b}{2a} = \pm \frac{\sqrt{\Delta}}{2a}
  \]
  Subtracting \(\frac{b}{2a}\) from both sides gives the famous quadratic formula:
  \[
  x = \frac{-b \pm \sqrt{\Delta}}{2a}
  \]
\end{itemize}

You could now go back and verify this formula in the previous example.

\newpage
\phantomsection
\section*{Conclusion}
\addcontentsline{toc}{section}{Conclusion}

\end{document}
